% ID: SEL-D0033-OB-E02
% ESTADO: validada. S034, cotejada con examen y criterios.
\textbf{a)} Cada tramo es continuo en su interior.
En $40$, el límite izquierdo es $-64+96=32$ y el valor
y límite derecho son $40-320/40=32$. Es continua en todo
$[0,50]$. Las derivadas son $2.4-0.08t$ y $320/t^2$.
En el empalme las pendientes son $-0.8$ y $0.2$, así que
no es derivable allí. En los extremos hay derivadas laterales.

\textbf{b)} La primera derivada se anula en $30$. El beneficio
crece en $[0,30]$, decrece en $[30,40]$ y crece en $[40,50]$.
Comparamos $B(30)=36$ con $B(50)=33.6$ y $B(0)=0$.
El máximo es de $\boxed{36000}$ euros en el año $\boxed{30}$.

\textbf{c)} El beneficio sube hasta el año 30, baja a 32000 euros
en el 40 y remonta hasta 33600 en el 50, por debajo del máximo.
\begin{center}
\setlength{\unitlength}{7mm}
\begin{picture}(14,10)(-1,-1)
\qbezier(0,0)(5.5,0)(11,0)\put(11.2,0){$t$}
\qbezier(0,0)(0,4)(0,8)\put(0,8.2){$B$}
\qbezier(0,0)(4,9.6)(8,6.4)
\qbezier(8,6.4)(9,6.6)(10,6.72)
\put(6,7.4){\makebox(0,0)[b]{$(30,36)$}}
\put(8,6.2){\makebox(0,0)[t]{$(40,32)$}}
\put(10.2,6.7){$(50,33.6)$}
\end{picture}
\end{center}
